% Page « L'essentiel » (chiffres : sorties/synthese_ensemble.txt et sorties des parties ; base 2025 le 17/11/2026 ;
% coupée le 17/11/2026, CONSIGNES_COUPE.md)
\section*{L'essentiel}

Cette note chiffre quatre réformes qui changent la manière de prélever sans changer ce que l'on prélève : chacune
rapporte autant que les impôts qu'elle remplace et laisse la répartition de la charge à peu près inchangée.

\begin{itemize}\setlength{\itemsep}{1pt}
\item \textbf{Fusionner l'impôt sur le revenu, la CSG et la CRDS} en un impôt individuel, le même pour tous les revenus,
avec un barème sans marches, une franchise sur les premiers euros d'activité et une franchise par enfant. Le premier
salaire au Smic serait prélevé à 7,8~\% au lieu de 9,5~\%.
\item \textbf{Taxer le sol plutôt que les ventes} : diviser par deux les droits de mutation et les remplacer par un
taux national d'environ 0,12~\% sur la valeur de tous les terrains privés.
\item \textbf{Déduire immédiatement l'investissement en équipements} plutôt que baisser le taux de l'impôt sur les
sociétés, en finançant la mesure par la suppression de niches.
\item \textbf{Garder le prélèvement forfaitaire unique, mais fermer la sortie sans impôt} : l'héritier ou le
donataire reprend le prix d'acquisition des titres reçus.
\end{itemize}

\textbf{Ce que l'on peut en attendre.} Environ 6,8~Md€ par an de perte sèche évitée, 0,23~\% du PIB (de 2,4 à
15,1~Md€ selon les hypothèses), dont 4,7 pour la fusion et 2,1 pour l'immobilier. Ce chiffre ne compte que les
réactions réelles (travail, emploi, ventes, investissement), pas l'optimisation fiscale. La transition se finance par
une avance qui culmine à environ 32~Md€ en 2032 et s'éteint vers 2040. Les montants sont ceux des prélèvements de
2025 ; la répartition de l'impôt sur le revenu entre foyers, celle des revenus de 2024.

\textbf{Qui paie.} Les actifs aux revenus modestes et moyens y gagnent : un jeune couple avec deux enfants qui achète
à Lyon gagnerait environ 3\,000~€ par an, droits de mutation compris. Les retraités paieraient davantage, parce qu'à
revenu égal ils seraient imposés comme les actifs : 21,5~Md€ de plus par an au total, 14,3 avec une franchise
d'activité réduite de moitié ; un retraité modeste propriétaire à Limoges, environ 1\,300~€ de plus par an.

\textbf{Ce qui peut faire changer la conclusion.} Le gain de la fusion suppose que les salariés réagissent aujourd'hui
à la CSG comme à l'impôt sur le revenu ; aucune étude ne le mesure, et s'ils la perçoivent comme une cotisation, le gain
total tombe à zéro, voire un peu en dessous. Celui de l'immobilier suppose que les ventes réagissent durablement aux droits. Enfin, la fin de la
contribution exceptionnelle sur les bénéfices des grandes entreprises, prévue par le droit en vigueur, retirerait
7~Md€ par an de recette : un choix distinct de nos réformes, à assumer ou à compenser.

\textbf{Ce que le décideur doit trancher} : l'ampleur de la franchise d'activité, c'est-à-dire ce que l'on demande
aux retraités pour encourager la reprise d'emploi ; l'impôt individuel ou l'imposition commune des couples, qui
réduit le risque constitutionnel ; le rythme de la transition.
